% Generated draft: requires mathematical and source review.
\documentclass[11pt]{article}
\usepackage[margin=25mm]{geometry}
\usepackage{fontspec}
\setmainfont{Latin Modern Roman}
\usepackage{amsmath,amssymb,mathtools,unicode-math}
\setmathfont{Latin Modern Math}
\usepackage{enumitem,needspace,xurl,hyperref}
\setlength{\parindent}{0pt}
\setlength{\parskip}{5pt}
\setlength{\emergencystretch}{3em}
\newcommand{\N}{\mathbb{N}}
\newcommand{\Z}{\mathbb{Z}}
\newcommand{\Q}{\mathbb{Q}}
\newcommand{\R}{\mathbb{R}}
\newcommand{\C}{\mathbb{C}}
\newcommand{\F}{\mathbb{F}}
\newcommand{\A}{\mathbb{A}}
\newcommand{\PP}{\mathbb{P}}
\newcommand{\E}{\mathbb{E}}
\newcommand{\eps}{\varepsilon}
\newcommand{\dd}{\,\mathrm{d}}
\begin{document}
\section*{UnsolvedMath ID: 10400073}
\textbf{Source problem code:} catalogue entry
\subsection{Definitions and mathematical statement}
For a surface \(\Sigma\), consider links in \(\Sigma\times[0,1]\) and finite-type invariants defined by alternating resolution of transverse double points. Let \(\mathcal A_\Sigma\) be the surface-diagram algebra used in the source; its precise diagram, relation, coefficient and completion conventions are not supplied.

T. Kohno asks for an \emph{explicit universal finite-type invariant} of these links with values in \(\mathcal A_\Sigma\). In the standard graded formulation, its degree-\(d\) truncation should recover every invariant of order at most \(d\) by linear evaluation, compatibly with the symbol/diagram correspondence. If the invariant is a formal series, its target is the corresponding degree completion; whether the source's notation \(\mathcal A_\Sigma\) already incorporates completion must be checked. The request is for an explicit construction, not just an abstract existence or a family of nonuniversal scalar invariants.

The source reports genus zero solved by the Kontsevich integral, a higher-genus construction suggestion by Deligne at Oberwolfach in 1995, and the punctured-surface case solved by Andersen, Mattes and Reshetikhin. These are supplied historical progress statements; they do not specify a completed construction for the remaining surface classes.

The appended context fixes a simple Lie group \(G\) and the moduli space \(M_G(\Sigma)\) of flat \(G\)-connections. It states that its smooth-function algebra \(\mathcal C(M_G(\Sigma))\) is Poisson, as is \(\mathcal A_\Sigma\), and that coloring each diagram component by a representation of \(G\) yields a natural Poisson algebra homomorphism
\[
 \tau:\mathcal A_\Sigma\longrightarrow\mathcal C(M_G(\Sigma)),\qquad
 \tau(\{a,b\})=\{\tau(a),\tau(b)\}.
\]
This is retained as context for the next problem announced but not supplied. Smooth-locus or singular-moduli conventions and the symplectic/Poisson hypotheses are not stated; a globally smooth symplectic moduli space is not assumed silently. The excerpt also leaves the topology/orientation of \(\Sigma\) and link framing conventions unresolved.
\subsection{Short English statement}
Explicitly construct a universal finite-type invariant for links in a thickened surface, valued in the source's surface-diagram algebra. Preserve the reported genus-zero and punctured-surface progress and its Poisson/flat-connection context, with missing conventions marked.
\subsection{Sources}
\begin{itemize}
\item Local UnsolvedMath registry, permanent ID 10400073, source code catalogue entry; statement attribution T. Kohno; proposer field and proposal year absent; update 2026-08-06.
\item Ohtsuki, \emph{Problems on invariants of knots and 3-manifolds} (2002), Problem 3.23, PDF page 71: \url{https://msp.org/gtm/2002/04/gtm-2002-04-024s.pdf}. Registry access date 2026-07-29; PDF font-structure extraction. Provenance distinguishes 231 numbered bold headings. Present status and redistribution terms require review.
\item Supplied historical attributions: Kontsevich integral in genus zero; Deligne's Oberwolfach 1995 suggestion; Andersen--Mattes--Reshetikhin for punctured surfaces. The Poisson-algebra statement cites source reference [8], whose bibliography entry is absent. No exact additional resolving reference is supplied.
\end{itemize}
\end{document}
